% !TeX root = ../main.tex
\section{Experimental Setup}
\label{sec:experimental-setup}

We evaluate ResNet-18, ConvNeXt-Tiny, and Tiny-ViT-5M backbones across CIFAR-10, CIFAR-100, and Flowers-102. This set provides a progression from entry-level convolutional validation to a compact transformer deployment target. The digital FP32 baselines are summarized in Table~\ref{tab:digital_baselines}. Models for CIFAR-10/100 were trained from scratch, while Flowers-102 used ImageNet-pretrained weights.

\begin{table}[htbp]
\centering
\caption{\textbf{Evaluated backbones and digital (FP32) accuracy baselines.} These baselines define the ideal software performance before analog non-ideality injection.}
\label{tab:digital_baselines}
\begin{tabular}{lccc}
\toprule
\textbf{Dataset} & \textbf{ResNet-18} & \textbf{ConvNeXt-Tiny} & \textbf{Tiny-ViT-5M} \\
\midrule
CIFAR-10 & 95.46\% & 90.74\% & 98.06\% \\
CIFAR-100 & 78.64\% & 64.12\% & 86.94\% \\
Flowers-102 & n.a. & 33.22\% & 97.97\% \\
\bottomrule
\end{tabular}
\end{table}

The main text then focuses on Tiny-ViT-5M under two regimes: an extreme low-precision 4-bit ablation used to diagnose cross-instance transfer, and a realistic phase-change memory (PCM) precision ladder spanning 4/6/8-bit operation. Detailed optimization settings, Monte Carlo evaluation protocols, and auxiliary architectural experiments, including ConvNeXt retention and proportional-noise studies, are summarized in the Supplementary Information.

Parameter provenance is summarized in the Supplementary Information together with the optimization settings and evaluation conditions used for the reported runs. The study targets reproducible reruns through archived configurations, logs, and model lineage. Because the framework operates at the simulation level, device parameters are literature-derived or proxy-calibrated rather than extracted from fabricated devices under direct measurement. The quantitative results should therefore be interpreted as comparative risk rankings and deployment frontiers for the tested profiles, not as absolute hardware predictions.
